\documentclass[a4paper]{article}

%use the english line for english reports
%usepackage[english]{babel}
\usepackage[portuguese]{babel}
\usepackage[utf8]{inputenc}
\usepackage{indentfirst}
\usepackage{graphicx}
\usepackage{verbatim}
\usepackage{fancyhdr}
\usepackage{listings}
\usepackage{color}

\definecolor{dkgreen}{rgb}{0,0.6,0}
\definecolor{gray}{rgb}{0.5,0.5,0.5}
\definecolor{mauve}{rgb}{0.58,0,0.82}

\lstset{frame=tb,
  language=Prolog,
  aboveskip=3mm,
  belowskip=3mm,
  showstringspaces=false,
  columns=flexible,
  basicstyle={\small\ttfamily},
  numbers=none,
  numberstyle=\tiny\color{gray},
  keywordstyle=\color{blue},
  commentstyle=\color{dkgreen},
  stringstyle=\color{mauve},
  breaklines=true,
  breakatwhitespace=true,
  tabsize=3
  }

\begin{document}


\setlength{\textwidth}{16cm}
\setlength{\textheight}{22cm}

\title{\Huge\textbf{SKYSCRAPER}\linebreak\linebreak\linebreak
\Large\textbf{Relatório}\linebreak\linebreak
\linebreak\linebreak
\includegraphics[scale=0.1]{images/feup-logo.png}\linebreak\linebreak
\linebreak\linebreak
\Large{Mestrado Integrado em Engenharia Informática e Computação} \linebreak\linebreak
\Large{Programação em Lógica}\linebreak
}

\author{\textbf{Grupo Skyscraper\_5:}\\
\linebreak\\
André Cruz - 201503776 \\
Edgar Carneiro - 201503748 \\
\linebreak\linebreak \\
 \\ Faculdade de Engenharia da Universidade do Porto \\ Rua Roberto Frias, s\/n, 4200-465 Porto, Portugal \linebreak\linebreak\linebreak
\linebreak\linebreak\vspace{1cm}}

\maketitle
\thispagestyle{empty}

%************************************************************************************************
%************************************************************************************************

\newpage

\section{Resumo}

Deve contextualizar e resumir o trabalho, salientando o objetivo, o método utilizado e fazendo referência aos principais resultados e à principal conclusão que esses resultados permitem obter.

\newpage

\tableofcontents

%************************************************************************************************
%************************************************************************************************

%*************************************************************************************************
%************************************************************************************************

\newpage

%%%%%%%%%%%%%%%%%%%%%%%%%%
\section{Introdução}

Descrição dos objetivos e motivação do trabalho, referência sucinta ao problema em análise (idealmente, referência a outros trabalhos sobre o mesmo problema e sua abordagem), e descrição sucinta da estrutura do resto do artigo.

\newpage

%%%%%%%%%%%%%%%%%%%%%%%%%%
\section{Descrição do Problema}

Descrever com detalhe o problema de otimização ou decisão em análise.

\newpage

%%%%%%%%%%%%%%%%%%%%%%%%%%
\section{Abordagem}

 Descrever a modelação do problema como um PSR, de acordo com as subsecções seguintes:

\subsection{Variáveis de Decisão}

Descrever as variáveis de decisão e os seus domínios.

\subsection{Restrições}

Descrever as restrições rígidas e flexíveis do problema e a sua implementação utilizando o SICStus Prolog.

\subsection{Função de Avaliação}

Descrever, quando for o caso, a forma de avaliar a solução obtida e a sua implementação utilizando o SICStus Prolog.

\subsection{Estratégia de Pesquisa}

Descrever a estratégia de etiquetagem (labeling) utilizada ou implementada, nomeadamente no que diz respeito à ordenação de variáveis e valores.

\newpage

%%%%%%%%%%%%%%%%%%%%%%%%%%

\section{Visualização da Solução}

Explicar os predicados que permitem visualizar a solução em modo de texto.

\newpage

%%%%%%%%%%%%%%%%%%%%%%%%%%
\section{Resultados}

Demonstrar exemplos de aplicação em instâncias do problema com diferentes complexidades e analisar os resultados obtidos. Devem ser utilizadas formas convenientes para apresentação dos resultados (tabelas e/ou gráficos).

\newpage

%%%%%%%%%%%%%%%%%%%%%%%%%%
\section{Conclusões e Trabalho Futuro }

Que conclusões retira deste projeto? O que mostram os resultados obtidos? Quais as vantagens e limitações da solução proposta? Como poderia melhorar o trabalho desenvolvido?

\clearpage
\addcontentsline{toc}{section}{Bibliografia}
\renewcommand\refname{Bibliografia}
\bibliography{myrefs}
\bibliographystyle{ieeetr}


\newpage
\appendix
\section{Anexo}

\huge\textbf{fabrik.pl}
\begin{lstlisting}
:- include('board.pl').
:- include('display.pl').
:- include('utils.pl').
:- include('menus.pl').
:- include('input.pl').
:- include('ai.pl').
:- include('test.pl').
:- use_module(library(random)).


fabrik:-
	mainMenuHandler.

%Game Initialization - set Workers and choose who starts
initGame(Player1, Player2) :-
	genRowColFacts,
	boardSize(N), !,
	createBoard(B0, N), printBoard(B0),
	setFirstWorker(Player1, black, B0, B1), printBoard(B1),
	setFirstWorker(Player2, black, B1, B2), printBoard(B2),
	chooseStartingPlayer(Player1, Side), printBoard(B2), !,
	gameLoop(Player1, Player2, Side, B2).
%If somehting failed on initGame, return to Play Menu
initGame(_, _):-
	playMenuHandler.

%game Loop for Player 1
gameLoop(Player1Function, Player2Function, black, Board) :-
	boardIsNotFull(Board),
	call(Player1Function, black, Board, NewBoard), printBoard(NewBoard),
	decideNextStep(Player1Function, Player2Function, black, NewBoard), !.
%Player 1 could not place a piece, so he lost
gameLoop(_Player1Function, _Player2Function, black, _Board) :- victory(white), !.

%game Loop for Player 2
gameLoop(Player1Function, Player2Function, white, Board) :-
	boardIsNotFull(Board),
	call(Player2Function, white, Board, NewBoard), printBoard(NewBoard),
	decideNextStep(Player1Function, Player2Function, white, NewBoard), !.
%Player 2 could not place a piece, so he lost
gameLoop(_Player1Function, _Player2Function, white, _Board) :- victory(black), !.


%Decide If someone has won with N in a row, or if the game should progress
decideNextStep(_Player1Function, _Player2Function, Side, Board) :-
	gameIsWon(Side, Board), !,
	victory(Side).
decideNextStep(Player1Function, Player2Function, Side, Board) :-
	changePlayer(Side, NewSide), !,
	gameLoop(Player1Function, Player2Function, NewSide, Board).

%End game with victory of Side
victory(Side):-
	destroyRowColFacts, !,
	wonMsg(Side),
	getEnter, !.

%Executes a human Play, getting worker input and piece input
humanPlay(Side, Board, NewBoard) :-
	workerUpdate(Side, Board, TempBoard),
	printBoard(TempBoard),
	isPiecePlayPossible(TempBoard), !,
	pieceInput(Side, Side, TempBoard, NewBoard), !.

%Setting the First Workers on the Board - Game beggining
setFirstWorker('humanPlay', Side, Board, NewBoard) :-
	pieceInput(worker, Side, Board, NewBoard).
%For AI Functions
setFirstWorker(_, _Side, Board, NewBoard) :-
        boardSize(Size),
        random(0, Size, Row), random(0, Size, Col),
        setPiece(worker, Row, Col, Board, NewBoard).

%Choose the Player that starts putting pieces on the board
chooseStartingPlayer('humanPlay', Side) :-
        getFirstPlayer(Side), !.
chooseStartingPlayer(_, white). % always white - For AI functions 
\end{lstlisting}
\newpage

\huge\textbf{board.pl}
\begin{lstlisting}
:- use_module(library(lists)).
:- use_module(library(clpfd)).
:- use_module(library(sets)).

% Generate a board predicate with N x N empty spaces
createBoard(Board, N) :-
        createBoard(Board, N , 0).

createBoard(_, N, N).
createBoard([FirstRow | OtherRows], N, Lines) :-
        Lines1 is (Lines + 1),
        createBoardLine(FirstRow, N),
        createBoard(OtherRows, N, Lines1).
        
createBoardLine(_, 0).
createBoardLine([FirstEle | OtherEle], N) :-
        FirstEle = none,
        N1 is (N - 1),
        createBoardLine(OtherEle, N1).


% Access the element in the [Row,Col] position of the given board
getElement(Board, Row, Col, Element):-
        nth0(Row, Board, RowLine, _),
        nth0(Col, RowLine, Element, _), !.


%% Validation Predicates
% Worker can be played if [Row,Col]=none
isValidPlay(worker, Row, Col, Board) :- !,
        getElement(Board, Row, Col, none).

% No conditions if there is no piece
isValidPlay(none, _, _, _) :- !.

% White/Black pieces can be played if [Row,Col] is in intersection of Workers' lines of sight
isValidPlay(_, Row, Col, Board) :-
        getElement(Board, Row, Col, none),
        isIntersection(Board, Row, Col).

% Fetches the position of both workers and checks if [Row,Col] is in their lines of sight
isIntersection(Board, Row, Col) :-
        findBothWorkers(Board, Row1, Col1, Row2, Col2),
        getIntersections(Board, Row1, Col1, Row2, Col2, Positions),
        member([Row, Col], Positions).

% Gets the intersections of the workers' lines of sight
getIntersections(Board, Row1, Col1, Row2, Col2, Positions) :-
        positionsInSight(Board, Row1, Col1, Pos1),
        positionsInSight(Board, Row2, Col2, Pos2),
        intersection(Pos1, Pos2, Positions).

% Position is in Board and is empty ?
isValidPosition(Board, Row, Col) :-
        boardSize(N),
        Row >= 0, Col >= 0,
        Row =< N, Col =< N,
        getElement(Board, Row, Col, none).

% Possible values for coordinates' change
coordinateChange(0).
coordinateChange(1).
coordinateChange(-1).

rowColChange(RowChange, ColChange) :-
        coordinateChange(RowChange),
        coordinateChange(ColChange),
        \+ (RowChange = 0, ColChange = 0).

% Spreads outwards from the worker's position and stops on end of board or when a piece blocks the line of sight
% Returns all positions of the worker's lines of sight
positionsInSight(Board, Row, Col, Positions) :-
        findall(PartialPositions, (rowColChange(RChange,CChange), lineOfSight(Board, Row, Col, RChange, CChange, PartialPositions)), ListOfLists),
        append(ListOfLists, Positions).

lineOfSight(Board, Row, Col, RowChange, ColChange, Positions) :-
        NewRow is Row + RowChange, NewCol is Col + ColChange,
        isValidPosition(Board, NewRow, NewCol), !,
        lineOfSight(Board, NewRow, NewCol, RowChange, ColChange, OtherPositions),
        append([[NewRow, NewCol]], OtherPositions, Positions).
lineOfSight(_Board, _Row, _Col, _RowChange, _ColChange, []) :- !.

                                                                   
% Set piece on board
% Sets the piece of the given type on the given position, on the given Board
setPiece(Piece, Row, Col, Board, NewBoard) :-
        isValidPlay(Piece, Row, Col, Board),
        nth0(Row, Board, RowLine, TmpBoard),
        nth0(Col, RowLine, _, TmpRowLine),
        nth0(Col, NewRowLine, Piece, TmpRowLine),
        nth0(Row, NewBoard, NewRowLine, TmpBoard).


% Finds both workers' positions
findBothWorkers(Board, Row1, Col1, Row2, Col2) :-
        findWorker(Board, Row1, Col1),
        findWorker(Board, Row2, Col2),
        \+ (Row1 = Row2, Col1 = Col2), !. % Cut prevents backtracking over Row/Col permutations 

% Finds a worker's position
findWorker(Board, OutputRow, OutputCol) :-
        nth0(OutputRow, Board, TmpRow),
        nth0(OutputCol, TmpRow, worker).

% Tests if board has any empty space
boardIsNotFull(Board) :-
        nth0(_, Board, TmpRow),
        nth0(_, TmpRow, none).

% Side has won ?
gameIsWon(PieceSide, Board):-
        checkHorizontalWin(PieceSide, Board).
gameIsWon(PieceSide, Board):-
        checkVerticalWin(PieceSide, Board).
gameIsWon(PieceSide, Board):-
        checkDiagonalWin(PieceSide, Board).

% Check Horizontal Win for side 'Side'
checkHorizontalWin(Side, [FirstRow | _RestOfBoard]) :-
        checkRowWin(Side, FirstRow, 0).
checkHorizontalWin(Side, [FirstRow | RestOfBoard]) :-
        \+ checkRowWin(Side, FirstRow, 0),
        checkHorizontalWin(Side, RestOfBoard).
checkRowWin(_, _, N) :-
        winningStreakN(N), !.
checkRowWin(Side, [Side | RestOfRow], Count) :-
        NewCount is Count + 1,
        checkRowWin(Side, RestOfRow, NewCount).
checkRowWin(Side, [FirstEl | RestOfRow], _Count) :-
        Side \= FirstEl,
        checkRowWin(Side, RestOfRow, 0).


% Check Vertical Win for side 'Side'
checkVerticalWin(Side, Board) :-
        transpose(Board, TransposedBoard),
        checkHorizontalWin(Side, TransposedBoard).

%Check Diagonal Win for side 'Side'
checkDiagonalWin(Side, Board):-
        boardSize(BoardSize),
        checkDiagWinAux(Side, Board, 0, BoardSize).

%One quarter of the diagonal lines: left-right, top-down
checkDiagWinAux(Side, Board, Col, _BoardSize):-
        checkDiagLineWin(Side, Board, 0, Col, 1, 1, 0).
%One quarter of the diagonal lines: right-left, top-down
checkDiagWinAux(Side, Board, Col, _BoardSize):-
        checkDiagLineWin(Side, Board, 0, Col, 1, -1, 0).
%One quarter of the diagonal lines: left-right, down-top
checkDiagWinAux(Side, Board, Col, BoardSize):-
        BottomRow is (BoardSize - 1),
        checkDiagLineWin(Side, Board, BottomRow, Col, -1, 1, 0).
%One quarter of the diagonal lines: right-left, down-top
checkDiagWinAux(Side, Board, Col, BoardSize):-
        BottomRow is (BoardSize - 1),
        checkDiagLineWin(Side, Board, BottomRow, Col, -1, -1, 0).
%Recursive Call
checkDiagWinAux(Side, Board, Col, BoardSize):-
        NewCol is (Col + 1),
        NewCol \= BoardSize,
        checkDiagWinAux(Side, Board, NewCol, BoardSize).

%Iterates through the diagonal line, starting at [Row, Col], with direction Vector [RowInc, ColInc], checking for Winning Streak
%Found N in-a-row, won game!
checkDiagLineWin(_Side, _Board, _Row, _Col, _RowInc, _ColInc, Count):-
        winningStreakN(Count), !.
checkDiagLineWin(Side, Board, Row, Col, RowInc, ColInc, Count):-
        getElement(Board, Row, Col, Side),
        NewRow is (Row + RowInc), NewCol is (Col + ColInc),
        NewCount is (Count + 1),
        checkDiagLineWin(Side, Board, NewRow, NewCol, RowInc, ColInc, NewCount).
checkDiagLineWin(Side, Board, Row, Col, RowInc, ColInc, _Count):-
        getElement(Board, Row, Col, _),
        NewRow is (Row + RowInc), NewCol is (Col + ColInc),
        checkDiagLineWin(Side, Board, NewRow, NewCol, RowInc, ColInc, 0).

%Moves a worker from the position [Row, Col] to the position [DestRow, DestCol], if possible
moveWorker(Board, Row, Col, DestRow, DestCol, UpdatedBoard):-
        getElement(Board, Row, Col, worker),
        setPiece(none, Row, Col, Board, TempBoard),
        setPiece(worker, DestRow, DestCol, TempBoard, UpdatedBoard).

% Checks if any play is possible (with the worker already having been placed)
isPiecePlayPossible(Board) :-
        findBothWorkers(Board, R1, C1, R2, C2),
        getIntersections(Board, R1, C1, R2, C2, PossiblePlays),
        length(PossiblePlays, Length),
        Length > 0, !.
\end{lstlisting}
\newpage

\huge\textbf{ai.pl}
\begin{lstlisting}
:- use_module(library(random)).
:- dynamic(validCoord/1).

% Defense factor in AI evaluation function
defenseFactor(Value, NewValue):-
	NewValue is (Value * 1).

%Destroy the board facts generated for every game
destroyRowColFacts:-
	retractall(validCoord(_)). 

% Possible values for Row and Col positions
genRowColFacts:-
	boardSize(N),
	genRowColFactsAux(0, N), !.

genRowColFactsAux(Current, Current).
genRowColFactsAux(Current, BoardSize) :-
	asserta(validCoord(Current)),
	NewValue is (Current + 1),
	genRowColFactsAux(NewValue, BoardSize).

% Used to backtrace over the possible positions
validCoordinates(RowChange, ColChange) :-
        validCoord(RowChange),
        validCoord(ColChange).

% returns a List containing all the possible Boards by moving the workers
getAllWorkerPermutations(Board, PossibleBoards) :-
	findBothWorkers(Board, Row1, Col1, Row2, Col2),
	getWorkerBoards(Board, Row1, Col1, PossibleBoards1),
	getWorkerBoards(Board, Row2, Col2, PossibleBoards2),
	union(PossibleBoards1, PossibleBoards2, PossibleBoards). % No duplicates

% All the boards obtained by moving a worker through a Board
getWorkerBoards(Board, Row, Col, ListOfBoards) :-
	findall(PossibleBoard, (validCoordinates(NewRow, NewCol), moveWorker(Board, Row, Col, NewRow, NewCol, PossibleBoard)), ListOfBoards), !.

% Get all the boards where a piece can go, given a List of Boards with different worker positions
getPieceBoards(Side, WorkerBoards, PossibleBoards) :-
	getPieceBoardsAux(Side, WorkerBoards, [], PossibleBoards).

getPieceBoardsAux(_, [], PossibleBoards, PossibleBoards) :- !.
getPieceBoardsAux(Side, [WorkerBoard | OtherBoards], SoFarBoards, PossibleBoards) :-
	findBothWorkers(WorkerBoard, Row1, Col1, Row2, Col2),
	getIntersections(WorkerBoard, Row1, Col1, Row2, Col2, IntersectionsList),
	genNewBoards(Side, WorkerBoard, IntersectionsList, [], NewBoards),
	append(NewBoards, SoFarBoards, UpdatedBoards), !,
	getPieceBoardsAux(Side, OtherBoards, UpdatedBoards, PossibleBoards).

% Generates all the boards associated to a certain board with fix workers.
% Boards with all the different places where the piece can be played
genNewBoards(_, _, [], AllBoards, AllBoards) :- !.
genNewBoards(Side, Board, [Intersec | OtherIntersec], FoundBoards, AllBoards) :-
	Intersec = [Row, Col],
	setPiece(Side, Row, Col, Board, TempBoard),
	append([TempBoard], FoundBoards, UpdatedBoards), !,
	genNewBoards(Side, Board, OtherIntersec, UpdatedBoards, AllBoards).

% Returns all the possible resulting boards, taking into account the move worker play and the set piece play
getPossibleBoards(Side, Board, PossibleBoards) :-
	getAllWorkerPermutations(Board, WorkerBoards), !,
	getPieceBoards(Side, WorkerBoards, PossibleBoards),
	length(PossibleBoards, Size),
	Size > 0. %If no board is created, game is Over


% Evaluates a board and returns the correspondent Value
evaluateBoard(Side, Board, BoardValue) :-
	horizontalEvaluation(Side, Board, HorizontalValue),
	verticalEvaluation(Side, Board, VerticalValue),
	diagonalEvaluation(Side, Board, DiagonalValue),
	defensiveEvaluation(Side, Board, DefensiveValue),
	BoardValue is (HorizontalValue + VerticalValue + DiagonalValue + DefensiveValue).

% Makes an Horizontal Evaluation of the given board
horizontalEvaluation(Side, Board, Value) :-
	horizontalEvaluationAux(Side, Board, 0 , Value), !.
horizontalEvaluationAux(_, [], FinalValue, FinalValue) :- !.
horizontalEvaluationAux(Side, [FirstRow | RestOfBoard], CurrentValue, FinalValue) :-
	horizontalRowEvaluation(Side, FirstRow, 0, 0, CurrentValue, LineFValue),
	horizontalEvaluationAux(Side, RestOfBoard, LineFValue, FinalValue).

horizontalRowEvaluation(_, [], _, _, LineFValue, LineFValue) :- !.
% Succession of Side Pieces
horizontalRowEvaluation(Side, [Side | OtherCols], Streak, _EnemyStreak, CurrentValue, LineFValue) :-
	NewStreak is (Streak + 1),
	NewValue is (CurrentValue + (NewStreak * NewStreak * NewStreak)),
	% neighborEnemyStreaks
	horizontalRowEvaluation(Side, OtherCols, NewStreak, 0, NewValue, LineFValue).
% Succession of Enemy Pieces
horizontalRowEvaluation(Side, [Col | OtherCols], _Streak, EnemyStreak, CurrentValue, LineFValue) :-
	changePlayer(Side, Enemy),
	Col = Enemy,
	NewEnemyStreak is (EnemyStreak + 1),
	NewValue is (CurrentValue - (NewEnemyStreak * NewEnemyStreak * NewEnemyStreak)),
	horizontalRowEvaluation(Side, OtherCols, 0, NewEnemyStreak, NewValue, LineFValue).
% When nor white nor black
horizontalRowEvaluation(Side, [_Col | OtherCols], _, _, CurrentValue, Value) :-
	horizontalRowEvaluation(Side, OtherCols, 0, 0, CurrentValue, Value).

% Makes a vertical Evaluation of the given board
verticalEvaluation(Side, Board, Value) :-
	transpose(Board, TransposedBoard),
	horizontalRowEvaluation(Side, TransposedBoard, 0, 0, 0, Value).

% Makes a diagonal Evaluation of the given board
diagonalEvaluation(Side, Board, Value) :-
	boardSize(BoardSize),
	diagonalEvaluationAux(Side, Board, 0, BoardSize, 0, Value).
diagonalEvaluationAux(_Side, _Board, BoardSize, BoardSize, FinalValue, FinalValue) :- !.
diagonalEvaluationAux(Side, Board, Col, BoardSize, CurrentValue, Value) :-
	BottomRow is (BoardSize - 1),
	ColWithoutMainDiag is (Col - 1),
	% One quarter of the diagonal lines: left-right, top-down
	diagonalLine(Side, Board, 0, Col, 1, 1, 0, 0, CurrentValue, DiagLine1Value),
	% One quarter of the diagonal lines: right-left, top-down
	diagonalLine(Side, Board, 0, ColWithoutMainDiag, 1, -1, 0, 0, DiagLine1Value, DiagLine2Value),
	% One quarter of the diagonal lines: left-right, down-top
	diagonalLine(Side, Board, BottomRow, Col, -1, 1, 0, 0, DiagLine2Value, DiagLine3Value),
	% One quarter of the diagonal lines: right-left, down-top
	diagonalLine(Side, Board, BottomRow, ColWithoutMainDiag, -1, -1, 0, 0, DiagLine3Value, DiagLine4Value),
	NewCol is (Col + 1),
	diagonalEvaluationAux(Side, Board, NewCol, BoardSize, DiagLine4Value, Value).

% Iterates thorugh the diagonal line, starting at [Row, Col] with the direction Vector [RowInc, ColInc], updating the line value
% Succession of Side Pieces
diagonalLine(Side, Board, Row, Col, RowInc, ColInc, Streak, _EnemyStreak, CurrentValue, FinalValue) :-
	getElement(Board, Row, Col, Side),
	NewRow is (Row + RowInc), NewCol is (Col + ColInc),
	NewStreak is (Streak + 1),
	NewValue is (CurrentValue + (NewStreak * NewStreak * NewStreak)),
	diagonalLine(Side, Board, NewRow, NewCol, RowInc, ColInc, NewStreak, 0, NewValue, FinalValue).
% Sucession of Enemy Pieces
diagonalLine(Side, Board, Row, Col, RowInc, ColInc, _Streak, EnemyStreak, CurrentValue, FinalValue) :-
	changePlayer(Side, Enemy),
	getElement(Board, Row, Col, Enemy),
	NewRow is (Row + RowInc), NewCol is (Col + ColInc),
	NewEStreak is (EnemyStreak + 1),
	NewValue is (CurrentValue - (NewEStreak * NewEStreak * NewEStreak)),
	diagonalLine(Side, Board, NewRow, NewCol, RowInc, ColInc, 0, NewEStreak, NewValue, FinalValue).
% When nor white nor black
diagonalLine(Side, Board, Row, Col, RowInc, ColInc, _Streak, _EnemyStreak, CurrentValue, FinalValue) :-
	getElement(Board, Row, Col, _),
	NewRow is (Row + RowInc), NewCol is (Col + ColInc),
	diagonalLine(Side, Board, NewRow, NewCol, RowInc, ColInc, 0, 0, CurrentValue, FinalValue).
% It only fails again if [Row, Col] out of board
diagonalLine(_Side, _Board, _Row, _Col, _RowInc, _ColInc, _Streak, _EnemyStreak, FinalValue, FinalValue) :- !.


defensiveEvaluation(Side, Board, Value) :-
	defensiveEvaluationRec(Side, Board, 0, 0, 0, Value).
% Piece is of type Side
defensiveEvaluationRec(Side, Board, Row, Col, CurrentValue, Value) :-
	getElement(Board, Row, Col, Side),
	updateDefenseValue(Side, Board, Row, Col, CurrentValue, UpdatedValue),
	getNextPosition(Row, Col, NewRow, NewCol),
	defensiveEvaluationRec(Side, Board, NewRow, NewCol, UpdatedValue, Value).
% Piece is not of type Side
defensiveEvaluationRec(Side, Board, Row, Col, CurrentValue, Value) :-
	getElement(Board, Row, Col, _),
	getNextPosition(Row, Col, NewRow, NewCol),
	defensiveEvaluationRec(Side, Board, NewRow, NewCol, CurrentValue, Value).
% Finished, accessed non-existent element - Finished Board
defensiveEvaluationRec(_, _, _, _, FinalValue, FinalValue) :- !.

% Evaluates Enemy Streaks near the given position
updateDefenseValue(Side, Board, Row, Col, CurrentValue, UpdatedValue) :-
	findall(LineValue, (validCoordinates(RChange,CChange), lineDefenseValue(Side, Board, Row, Col, RChange, CChange, 0, LineValue)), ListOfValues),
	append(ListOfValues, ListValues),
	sum_list(ListValues, Sum),
	UpdatedValue is (CurrentValue + Sum).

lineDefenseValue(Side, Board, Row, Col, RowChange, ColChange, Streak, Value) :-
	NewRow is Row + RowChange, NewCol is Col + ColChange,
	changePlayer(Side, Enemy),
        getElement(Board, NewRow, NewCol, Enemy), !,
        NewStreak is (Streak + 1),
        lineDefenseValue(Side, Board, NewRow, NewCol, RowChange, ColChange, NewStreak, OtherValues),
        TmpValue is (NewStreak * NewStreak * NewStreak),
        defenseFactor(TmpValue, PosValue),
        append([PosValue], OtherValues, Value).
lineDefenseValue(_Side, _Board, _Row, _Col, _RowChange, _ColChange, _Streak, _Value) :- !. % Needed?

% For-each function, TODO substitute if similar standard function is found
evaluateAllBoards(Side, [FirstBoard | OtherBoards], Result) :-
	evaluateBoard(Side, FirstBoard, FirstResult),
	evaluateAllBoards(Side, OtherBoards, TmpResult),
	append([FirstResult-FirstBoard], TmpResult, Result).
evaluateAllBoards(_, [], []).

getGreedyPlay(Side, CurrentBoard, Play) :-
	getPossibleBoards(Side, CurrentBoard, PossibleBoards),
	evaluateAllBoards(Side, PossibleBoards, GradedBoards),
	keysort(GradedBoards, Plays),
	last(Plays, _Grade-Play).

getRandomPlay(Side, CurrentBoard, Play) :-
	getPossibleBoards(Side, CurrentBoard, PossibleBoards),
	length(PossibleBoards, Len),
	Len1 is Len - 1,
	random(0, Len1, Idx),
	nth0(Idx, PossibleBoards, Play).
\end{lstlisting}
\newpage

\huge\textbf{utils.pl}
\begin{lstlisting}
%% CONSTANTS

boardSize(9) :- !.
%boardSize(11) :- !.

winningStreakN(5).

changePlayer(black, white).
changePlayer(white, black).

%Gets the next Position in a board
getNextPosition(Row, Col, Row, NCol):-
	NCol is (Col + 1),
	boardSize(Size),
	NCol < Size.

getNextPosition(Row, _Col, NRow, 0):-
	NRow is (Row + 1).

%Returns the sum of all elements of a List of Values
sum_list(List, Sum):-
	sum_listAux(List, 0, Sum).
sum_listAux([Head | ResList], ItSum, Sum):-
	NewSum is (ItSum + Head),
	sum_listAux(ResList, NewSum, Sum).
sum_listAux([], Sum, Sum).

%Prints all boards in an array of boards
printAllBoards([]).
printAllBoards([Board | NextBoards]):-
	boardSize(Size),
	printBoard(Board, Size),
	printAllBoards(NextBoards).

%Concatenates a list of of integers into an integer
concat_numbers(IntList, Int):-
	reverse(IntList, RevList),
	concat_numbersAux(RevList, 1, Int), !.
concat_numbersAux([], _, 0).
concat_numbersAux([FirstInt | OtherInts], TenMulti, Int):-
	NewTenM is (TenMulti * 10),
	concat_numbersAux(OtherInts, NewTenM, TmpInt),
	Int is (TenMulti * FirstInt + TmpInt).


clearConsole:-
	clearConsole(60).
clearConsole(0).
clearConsole(N) :-
	nl, 
	N1 is N-1, 
	clearConsole(N1).

printBoardEval([Val, Board]) :-
	printBoard(Board), nl,
	write('Value: '), write(Val), nl.
\end{lstlisting}
\newpage

\huge\textbf{display.pl}
\begin{lstlisting}
% Dictionary for Board Elements
translate(none, 32). %Empty Cell
translate(black, 79). %Dark Pieces
translate(white, 88). %White Pieces
translate(worker, 9608). %Red Workers

currentSideDisplay(Side):-
        write('   *** '), write(Side), write(' turn! ***'), nl, nl.

% General PrintBoard
printBoard(Board):-
        boardSize(N),
        printBoard(Board, N), !.

% Board Printing - arguments: Board and Board size
printBoard(Board, N):-
        clearConsole,
        write('  '), printHorizontalLabel(N, N),
        printBoard(Board, N, 1), !.

printBoard([], N, _):-
        printRowDivider(N), nl.

printBoard([Line | Board], N, CurrentL):-
        printRowDivider(N),
        printDesignRow(N),
        printVerticalLabel(CurrentL),
        put_code(9474),
        printLine(Line),
        printDesignRow(N),
        NewL is (CurrentL + 1),
        printBoard(Board, N, NewL).

printLine([]):- nl.
printLine([Head | Tail]) :-
        translate(Head, Code),
        write('   '),
        put_code(Code),
        write('   '), put_code(9474),
        printLine(Tail).

%  AESTHETICS

printRowDivider(N):-
        write('  '),
        put_code(9532),
        printRowDividerRec(N).

printRowDividerRec(0) :- nl.
printRowDividerRec(N) :-
        put_code(9472), put_code(9472), put_code(9472), put_code(9472),
        put_code(9472), put_code(9472), put_code(9472), put_code(9532),
        N1 is (N-1),
        printRowDividerRec(N1).

printDesignRow(N):-
        write('  '),
        put_code(9474),
        printDesignRowRec(N).

printDesignRowRec(0) :- nl.
printDesignRowRec(N) :-
        write('       '), put_code(9474),
        N1 is (N-1),
        printDesignRowRec(N1).

%Dictionary for Labels
getLabel( 0, 'A').
getLabel( 1, 'B').
getLabel( 2, 'C').
getLabel( 3, 'D').
getLabel( 4, 'E').
getLabel( 5, 'F').
getLabel( 6, 'G').
getLabel( 7, 'H').
getLabel( 8, 'I').
getLabel( 9, 'J').
getLabel(10, 'K').
getLabel(11, 'L').
getLabel(_,_):-
        write('Error: Unrecognized Label.'), nl,
        fail.

printHorizontalLabel(0, _):- nl.
printHorizontalLabel(N, Total):-
        Pos is (Total-N),
        getLabel(Pos, L),
        write('    '), write(L), write('   '),
        N1 is (N-1),
        printHorizontalLabel(N1, Total).        

printVerticalLabel(CurrentL):-
        CurrentL < 10,
        write(CurrentL),
        write(' ').

printVerticalLabel(CurrentL):-
        write(CurrentL).


printFabrikTitle:-
        clearConsole,
        write('   *********************************'), nl,
        write('   *    __  __   __   __           *'), nl,
        write('   *   |_  |__| |__) |__| | |__/   *'), nl,
        write('   *   |   |  | |__) | \\  | |  \\   *'), nl,
        write('   *                               *'), nl,
        write('   *********************************'), nl, nl.

wonMsg(Side):-
        wonArtMsg(Side).

wonArtMsg(white):-
        write('   *****************************************************'), nl,
        write('   *                    ___   __                       *'), nl,
        write('   *   \\    /  |__|  |   |   |_     \\    /  |   |\\ |   *'), nl,
        write('   *    \\/\\/   |  |  |   |   |__     \\/\\/   |   | \\|   *'), nl,
        write('   *                                                   *'), nl,
        write('   *****************************************************'), nl, nl.

wonArtMsg(black):-
        write('   *****************************************************'), nl,
        write('   *   __         __    __                             *'), nl,
        write('   *  |__)  |    |__|  /    |__/     \\    /  |   |\\ |  *'), nl,
        write('   *  |__)  |__  |  |  \\__  |  \\      \\/\\/   |   | \\|  *'), nl,
        write('   *                                                   *'), nl,
        write('   *****************************************************'), nl, nl.
\end{lstlisting}
\newpage

\huge\textbf{menus.pl}
\begin{lstlisting}
% Main Menu

mainMenu:-
	printFabrikTitle,
	write('\t 1. Play'), nl,
	write('\t 2. Rules'), nl,
	write('\t 3. Exit'), nl, nl,
	write('Choose an option:'), nl.

playMenu:-
	printFabrikTitle,
	write('\t 1. MultiPlayer'), nl,
	write('\t 2. SinglePlayer'), nl,
	write('\t 3. AI VS AI'), nl,
	write('\t 4. Back'), nl, nl,
	write('Choose an option:'), nl.

aiMenu:-
	printFabrikTitle,
        write('Choose the AI difficulty:'), nl,
        write('\t1. Easy AI'), nl,
        write('\t2. Smart AI'), nl, nl,
        write('Choose an option: '), nl.

rules:-
	printFabrikTitle,
	write('RULES: '), nl,
	write('    Beggining:'), nl,
	write('\tBlack starts by placing one of the workers on any space. Then White places the'), nl, 
	write('\tother worker on an arbitrary empty space. '), nl,
	write('\tBlack decides on who goes first. This player must place a stone of his color'), nl, 
	write('\taccording to the rules described below.'), nl, nl,
	write('    Objective:'), nl,
	write('\tPlayers win by creating a line of (at least) 5 stones in their color, orthogonally'), nl,
	write('\tor diagonally. Players lose the game immediately if they cannot place neither of'), nl,
	write('\tthe two workers in such a way that a new stone can be entered.'), nl, nl,
	write('    Play:'), nl,
	write('\tEach turn players may take - this is optional - one of the worker figures and'), nl,
	write('\tplace it on another empty space. After that, they must enter one of their stones'), nl,
	write('\ton an intersection point of the two workers\' lines of sight. These lines radiate'), nl,
	write('\tfrom a worker\'s position in orthogonal and diagonal directions arbitrarily far'), nl,
	write('\tover empty spaces.'), nl,
	write('\tNote: In the special case where the two workers are located on the same orthogonal'), nl,
	write('\tor diagonal line, all empty spaces between them are considered intersection points'), nl, nl,
	write('\t\t\t\t   Source: https://spielstein.com/games/fabrik/rules'), nl, nl,
	getEnter.

%User choose between Playing, Exiting or checking the rules
mainMenuHandler:-
	mainMenu,
	getInt(Choice),
	mainMenuChoice(Choice), !.
mainMenuChoice(1):-
	playMenuHandler, !.
mainMenuChoice(2):-
	rules,
	mainMenuHandler, !.
mainMenuChoice(3).
mainMenuChoice(_):-
	unknownInput,
	mainMenuHandler, !.

%User chooses between PvP, PvAI, AIvAI
playMenuHandler:-
	playMenu,
	getInt(Choice),
	playMenuChoice(Choice), !.
playMenuChoice(1):-
	initGame(humanPlay, humanPlay), !,
	mainMenuHandler, !.
playMenuChoice(2):-
	aiMenuHandler(AI),
        initGame(humanPlay, AI), !,
        mainMenuHandler, !.
playMenuChoice(3):-
	aiMenuHandler(AI1),
	aiMenuHandler(AI2),
        initGame(AI1, AI2), !,
        mainMenuHandler, !.
playMenuChoice(4):-
	mainMenuHandler, !.
playMenuChoice(_):-
	unknownInput,
	playMenuHandler, !.

%User chooses the AI Difficulty
aiMenuHandler(AI):-
        aiMenu,
        getInt(Choice),
        aiMenuChoice(Choice, AI), !.
aiMenuChoice(1, 'getRandomPlay').
aiMenuChoice(2, 'getGreedyPlay').
aiMenuChoice(_, AI):-
	unknownInput,
	aiMenuHandler(AI), !.
\end{lstlisting}
\newpage

\huge\textbf{input.pl}
\begin{lstlisting}
unknownInput:-
        write('Invalid option chosen.'), nl,
        getEnter.

%Asks the User form a integer.
%Reads the entire line
getInt(Input):-
        getIntCycle([], InputList),
        concat_numbers(InputList, Input).
getIntCycle(PrevInput, Input):-
        get_code(user_input, KbInput),
        evalCode(KbInput, PrevInput, Input).
evalCode(10, _, _):- !.
evalCode(Code, PrevInput, Input):-
        Elem is (Code - 48),
        getIntCycle(PrevInput, TmpInput),
        append([Elem], TmpInput, Input).

%Ask The User for a char
%Reads the entire line
getChar(Input):-
        get_char(user_input, KbInput),
        evalChar(KbInput, Input).
evalChar('\n', ''):- !.
evalChar(Code, Input):-
        getChar(TmpInput),
        atom_concat(Code, TmpInput, Input).

%'Press enter to continue' function
getEnter:-
        write('Press enter to continue.'), nl,
        getChar(_Input).

%Ask the User for a piece's row and column.
getPosition(PieceType, Row, Col):-
        getRow(PieceType, Row),
        getCol(PieceType, Col), !.

%Ask User for the Piece's row
getRow(PieceType, Row):-
        write('Choose '), write(PieceType), write('\'s row:'), nl,
        getInt(TempRow),
        Row is  (TempRow - 1), !.

%Ask User for the Piece's column
getCol(PieceType, Col):-
        write('Choose '), write(PieceType), write('\'s column:'), nl,
        getChar(ColLabel),
        getLabel(Col, ColLabel), !.

%User interface for placing a new piece on the board
pieceInput(PieceType, Side, Board, UpdatedBoard):-
        currentSideDisplay(Side),
        getPosition(PieceType, InputRow, InputCol),
        setPiece(PieceType, InputRow, InputCol, Board, UpdatedBoard).
pieceInput(PieceType, Side, Board, UpdatedBoard):-
        write('That play is not valid. Try again.'), nl, nl,
        pieceInput(PieceType, Side, Board, UpdatedBoard).

%User chooses who starts playing
getFirstPlayer(Side):-
        decidePlayerMsg,
        getInt(Choice),
        getFirstPlayerChoice(Choice, Side).
getFirstPlayer(Side):-
        write('Unrecognized choice. Try again.'), nl, nl,
        getFirstPlayer(Side).

getFirstPlayerChoice(1, white).
getFirstPlayerChoice(2, black).
decidePlayerMsg:-
        write('Black Player, decide who goes first:'), nl,
        write('\t1. White Player'), nl,
        write('\t2. BlackPlayer'), nl, nl,
        write('Choose an option:'), nl.


%User chooses if he wishes to update the worker position
workerUpdate(Side, Board, UpdatedBoard):-
        currentSideDisplay(Side), workerUpdateMsg,
        getInt(Choice),
        workerUpdateChoice(Choice, Side, Board, UpdatedBoard).

workerUpdateChoice(1, Side, Board, UpdatedBoard):-
        moveWorkerInput(Side, Board, UpdatedBoard), !.
workerUpdateChoice(2, _Side, Board, Board):- !.
workerUpdateChoice(_, Side, Board, UpdatedBoard):-
        unknownInput,
        workerUpdate(Side, Board, UpdatedBoard).

workerUpdateMsg:-
        write('Do you wish to move a worker?'), nl,
        write('\t1. Yes'), nl,
        write('\t2. No'), nl, nl,
        write('Choose an option:'), nl.

%Asks the User for worker's current and new positions and moves it
moveWorkerInput(Side, Board, UpdatedBoard):-
        currentSideDisplay(Side),
        write('\tWorker current position'), nl,
        getPosition(worker, CurrRow, CurrCol),
        write('\tWorker new position'), nl,
        getPosition(worker, DestRow, DestCol),
        moveWorker(Board, CurrRow, CurrCol, DestRow, DestCol, UpdatedBoard).
moveWorkerInput(Side, Board, UpdatedBoard):-
        write('That play is not valid. Try again.'), nl, 
        write('Help:'), nl,
        write('\t * To maintain the worker in the same place, keep the new position equal to the old one.'), nl,
        write('\t * The Worker must be moved to an empty cell.'), nl, nl,
        moveWorkerInput(Side, Board, UpdatedBoard).
\end{lstlisting}
\newpage

\end{document}
